\documentclass[a4paper,12pt]{report}
\usepackage{makecell}
\usepackage{placeins}
\usepackage[T5]{fontenc}
\usepackage[american]{babel}
\usepackage[margin=1in]{geometry}
\usepackage{tikz}
\usepackage{amssymb}
\usetikzlibrary{calc}
\usepackage{graphicx}
\usepackage{lmodern}
\usepackage{titlesec}
\titlespacing*{\chapter}{0pt}{0pt}{0pt}
\usepackage{fancyhdr}
\pagestyle{fancy}
\setlength{\headheight}{15pt}
\addtolength{\topmargin}{-2.5pt}
\fancyhead[L]{}
\fancyhead[R]{\textit{\MakeUppercase{\leftmark}}}
\usepackage{amsmath}
\usepackage{array, booktabs, multirow, longtable}
\usepackage{tabularx}
\usepackage{adjustbox}
\usepackage{newunicodechar}
\newunicodechar{≈}{\approx}
\usepackage{listings}
\usepackage{caption}
\usepackage{subcaption}
\usepackage{makecell}
\usepackage{float}
\usepackage{tocloft}
\usepackage[colorlinks=true, linkcolor=black, citecolor=black, urlcolor=black]{hyperref}
\usepackage{bookmark}
\usepackage{lipsum}
\usepackage{cite}
\usepackage{setspace}
\definecolor{codegreen}{rgb}{0,0.6,0}
\definecolor{codegray}{rgb}{0.5,0.5,0.5}
\definecolor{codepurple}{rgb}{0.58,0,0.82}
\definecolor{backcolour}{rgb}{0.95,0.95,0.92}
\lstdefinestyle{mystyle}{
    backgroundcolor=\color{backcolour},
    commentstyle=\color{codegreen},
    keywordstyle=\color{magenta},
    numberstyle=\tiny\color{codegray},
    stringstyle=\color{codepurple},
    basicstyle=\ttfamily\footnotesize,
    breakatwhitespace=false,
    breaklines=true,
    captionpos=b,
    keepspaces=true,
    numbers=left,
    numbersep=5pt,
    showlines=false,
    showspaces=false,
    showstringspaces=false,
    showtabs=false,
    tabsize=2
}
\lstset{style=mystyle, showlines=false}
% Glossaries are not used in this document, so disable the package to avoid warnings.
% \usepackage[acronym]{glossaries}
% \makeglossaries

\begin{document}

% ===== COVER PAGE =====

\thispagestyle{empty}
\begin{tikzpicture}[remember picture, overlay]
    \draw[line width=2pt]
        ($(current page.north east) - (0.5in, 0.5in)$);
\end{tikzpicture}
\begin{center}
    \adjustbox{max width=\textwidth}{\small \textbf{DEPARTMENT OF INFORMATION AND COMMUNICATION TECHNOLOGY}} \vspace{1cm}
%\includegraphics[width=0.7\textwidth]{images/usth.png} \vspace{1cm}
    % Move Thesis lower
    \vspace{0.7cm} \\
    {\Huge \textbf{Report}} \vspace{1.0cm}
    \rule{\textwidth}{0.75pt} \vspace{0.6cm}
    \LARGE \textbf{Threat Modelling in Large Language Models} \vspace{0.4cm}
    \rule{\textwidth}{0.75pt} \vspace{1.2cm}
    % Author
    {\Large \textbf{By Nguyen Huy Hai}}
\end{center}
\vfill
\begin{center}
    % Supervisor above date
    %{\large \textbf{Supervisor: Assoc. Prof. Tran Giang Son}} \\ \vspace{0.8cm}
    \textit{Ha Noi, August 2026}
\end{center}

\newpage


% ===== FRONT MATTER =====
\newpage
\pagenumbering{arabic}
\setcounter{page}{1}
\setlength{\parindent}{0pt}
\setlength{\parskip}{3pt}

\newpage
\tableofcontents

\newpage
\listoffigures
\addcontentsline{toc}{chapter}{List of Figures}

\newpage
\listoftables
\addcontentsline{toc}{chapter}{List of Tables}

\newpage
\chapter*{List of Abbreviations}
\addcontentsline{toc}{chapter}{List of Abbreviations}
\printglossary[type=\acronymtype]
{
\renewcommand{\arraystretch}{1.3}
\begin{longtable}{@{}lp{12cm}@{}}
\toprule
Term & Definition \\ 
\midrule
CARLA   & Car Learning to Act \\
\bottomrule
\end{longtable}
}
\markboth{List of Abbreviations}{}

\newpage
\chapter{Introduction}
% Why threat modelling is the first step 
Threat modeling is not an isolated preliminary step in this project, the organizing method that many component builds on.
\begin{itemize}
    \item \textbf{Intent Drift identification:} Threat modeling turns Intent Drift from a vague phenomenon into a defined taxonomy of attack vectors, a systematic account of how and where drift occurs.

    \item \textbf{Red Teaming:} That taxonomy defines the attack tree the Red Teaming AI searches, so the framework generates targeted adversarial test cases instead of relying on random prompt fuzzing.

    \item \textbf{Monitoring and defenses:} The same threat model specifies the deviation signatures, what drift looks like mechanistically, semantic integrity monitoring is built to detect, and that intent-aware defenses act on.

    \item \textbf{Multilingual/cross-cultural benchmarks:} Threat modeling identifies which attack patterns are genuinely language and culture-specific such as idiom, script-level tricks, culturally specific social-engineering. So the benchmark tests these directly rather than relying on English-language attacks that have simply been translated.
\end{itemize}

Together, these roles make threat modeling the shared reference model that unifies Red Teaming, monitoring, provenance, verification, and benchmarking into defense pipeline.

\chapter{Related Works}
% tell paper have done works related to

\textbf{ThreMoLIA — Threat Modeling of Large Language Model Integrated Applications}
\begin{itemize}
    \item \textbf{Problem/gap:} Threat modeling for LLM-integrated applications (LIAs) is manual and effort-intensive; no automated support exists.
    
    \item \textbf{Approach:} LLM + RAG-assisted framework that starts from a system Data Flow Diagram and combines MITRE ATLAS, OWASP Top 10 for LLMs, and STRIDE to generate/update a threat model, augmented with retrieval over existing threat models, security knowledge, and architecture docs.
    
    \item \textbf{Key result(s):} Proposes a "health score" concept (19 candidate metrics — e.g., threat coverage, residual risk, model complexity) to estimate threat-model completeness/quality; no concrete scoring formula defined yet.
    
    \item \textbf{Limitation:} Early-stage vision paper; evaluation limited to a single proof-of-concept case study on ChatGPT-3.5; no formula validated with practitioners; no industrial case studies.
    
    \item \textbf{Relevance to this work:} \textbf{[ASSUMPTION — confirm/edit]} ThreMoLIA's LLM+RAG-assisted, DFD-driven approach to generating a threat model is the closest prior work to the threat-modeling stage of this framework, but lacks the intent-drift-specific taxonomy and downstream automation (Red Teaming, monitoring, provenance) this work builds on top of a threat model.
\end{itemize}

\vspace{0.5em}
\textbf{Red Teaming Language Models with Language Models (Perez et al., EMNLP 2022)}
\begin{itemize}
    \item \textbf{Problem/gap:} Manual Red Teaming (humans hand-writing adversarial prompts) is costly and doesn't scale.
    
    \item \textbf{Approach:} A "red LM" automatically generates adversarial test inputs in a three-stage pipeline: (1) red LM generates adversarial test input $x$, (2) target LM produces reply $y$, (3) a classifier $r(x,y)$ flags harmful outputs. Compares four generation methods (zero-shot, stochastic few-shot, supervised learning, RL) on a diversity/difficulty trade-off.
    
    \item \textbf{Key result(s):} RL is most effective at eliciting harm but prone to mode collapse at aggressive settings (78\% of generated questions converged on "invisible"). Applied to a 280B-parameter chatbot, surfaced five harm categories: offensive replies, training-data leakage, generated contact info, distributional bias, and multi-turn compounding harms.
    
    \item \textbf{Limitation:} Effectiveness is bounded by the red LM and classifier, both of which inherit biases from their training data; complements rather than replaces manual Red Teaming.
    
    \item \textbf{Relevance to this work:} \textbf{[ASSUMPTION]} Using an LM to automate adversarial test generation, directly relevant to the AI-driven Red Teaming component of this framework; this work's threat model supplies that structure to target generation rather than searching unconstrained.
\end{itemize}

\vspace{0.5em}
\textbf{HarmBench — A Standardized Evaluation Framework for Automated Red Teaming and Robust Refusal}
\begin{itemize}
    \item \textbf{Problem/gap:} Prior automated Red Teaming evaluations lack breadth (most test 
$<100$ behaviors), comparability (e.g., completion length alone can swing ASR by up to 
30\%), and robust classifiers (many fail simple adversarial sanity checks, such as missing 
a model that initially refuses but then complies anyway).
    
    \item \textbf{Approach:} 510 curated harmful behaviors across 7 semantic categories and 4 functional categories, with an official validation/test split. Pipeline: generate test cases $\rightarrow$ generate completions (fixed 512 tokens, greedy decoding) $\rightarrow$ classify success via a fine-tuned classifier. Also introduces R2D2, an adversarial-training defense method. Recheck with (As an AI, i am sorry, …) and perplexity, not just ASR.
    
    \item \textbf{Key result(s):} Benchmarks 18 Red Teaming methods against 33 target LLMs/defenses; finds no attack or defense is uniformly effective, 
    and robustness does not scale with model size within a family — training data/procedure matters more. R2D2 achieves state-of-the-art robustness against 
    GCG while preserving general capability.
    
    \item \textbf{Limitation:} Standardization is scoped to the specific behavior set and classifiers provided; not stated to cover multilingual or 
    culturally-specific attack behaviors.
    
    \item \textbf{Relevance to this work:} \textbf{[ASSUMPTION]} HarmBench's emphasis on standardized, comparable evaluation is the model for this work's 
    own benchmarking approach, but its behavior set is not designed for the language and culture-specific attack patterns this work's multilingual benchmarks target.
\end{itemize}

\vspace{0.5em}
\textbf{MITRE ATLAS}
\begin{itemize}
    \item \textbf{Problem/gap:} Shared vocabulary existed for cataloging real-world adversarial attacks against AI/ML systems.
    
    \item \textbf{Approach:} A living knowledge base structured like MITRE ATT\&CK, organizing attacks into tactics (adversary objectives) and techniques 
    (specific methods), each with identifiers, descriptions, mitigations, and references to documented incidents.
    
    \item \textbf{Key result(s):} Continuously updated; recent releases expand coverage of generative-AI threats such as prompt injection and RAG-related attacks.
    
    \item \textbf{Limitation:} Not stated in the source material beyond general scope (curated/manual maintenance; framework-level rather than automated).    
    \item \textbf{Relevance to this work:} \textbf{[ASSUMPTION]} ATLAS serves as the technique-level checklist for this work's threat register, mapping 
    identified Intent Drift threats to ATLAS techniques gives a standardized, cross-referenceable vocabulary.
\end{itemize}

\textbf{Azure AI Content Safety}
Azure AI Content Safety is Microsoft's production content-moderation service, offered as 
part of Azure AI Foundry, for detecting harmful content in both user-generated and 
AI-generated text and images.

\begin{itemize}
    \item \textbf{Problem/gap:} Deployed LLM applications need a way to detect and grade 
    harmful content in real time, not just flag it as present or absent.
    \item \textbf{Approach:} An ensemble of classification models scores text (and images) 
    against four harm categories: Hate/Fairness, Sexual, Violence, and Self-Harm. Each 
    on an internal 0-7 severity scale. The API can also return a trimmed 4-level scale 
    (0, 2, 4, 6), where adjacent raw levels are collapsed together, letting applications 
    set per-category severity thresholds rather than a single pass/fail cutoff. The 
    broader service also includes Prompt Shields for prompt-injection detection and 
    groundedness detection for hallucination detection.
    \item \textbf{Key result(s):} Provides configurable, per-category severity thresholds 
    integrated natively into Azure OpenAI Service, so prompts and completions are filtered 
    automatically without additional application logic.
    \item \textbf{Limitation:} Categories are fixed to general content-harm types 
    (hate, sexual, violence, self-harm) rather than security-specific attack vectors; 
    it is not designed to classify Intent Drift or other reasoning-level security threats.
    \item \textbf{Relevance to this work:} [ASSUMPTION]
    Azure AI Content Safety's graded severity model (fine-grained internal scoring, 
    coarser reported levels) is a useful industry reference for grading the severity 
    of detected Intent Drift events, but its fixed harm categories do not map onto this 
    work's threat taxonomy, motivating a drift-specific severity scale.
\end{itemize}

\chapter{Methodology - Planning}
% tell proposed main steps here (just planning phase)
\section{Define the System Boundary}
Threat modelling is started with a DFD(Data Flow Diagram), then STRIDE is applied systematically to each component to identify: external entities, processes, 
databases, trust boundaries, and data flows. After that, the security bug hunt starts.

\section{Enumerate Threats with STRIDE}
STRIDE is a threat modeling framework developed by Microsoft that categorizes
security threats into six classes: Spoofing, Tampering, Repudiation,
Information Disclosure, Denial of Service, and Elevation of Privilege.
For LLM-integrated applications, these categories can be reinterpreted as
security questions that guide the identification of AI-specific threats, such
as prompt injection, hidden prompt leakage, and context exhaustion.

\begin{table}[ht]
\centering
\caption{Interpreting STRIDE for LLM Threat Modeling}
\label{tab:stride_llm}
\renewcommand{\arraystretch}{1.3}
\begin{tabular}{|p{2.2cm}|p{4.8cm}|p{6.5cm}|}
\hline
\textbf{Threat} &
\textbf{Question} &
\textbf{LLM Interpretation Example} \\
\hline
Spoofing &
Can someone pretend to be someone else? &
Can someone impersonate the LLM provider? \\
\hline
Tampering &
Can someone modify data or code? &
Can prompts be altered? \\
\hline
Repudiation &
Can someone deny performing an action? &
Can prompt history be deleted? \\
\hline
Information Disclosure &
Can confidential information leak? &
Can hidden prompts leak? \\
\hline
Denial of Service &
Can the system be made unavailable? &
Can huge prompts exhaust context? \\
\hline
Elevation of Privilege &
Can someone gain more permissions than intended? &
Can prompt injection bypass safety instructions? \\
\hline
\end{tabular}
\end{table}

Here there are some threat model matrix for LLM/LMM:
\begin{table*}[ht]
\centering
\caption{Example Threat Model Matrix for LLM/LMM}
\label{tab:threat_model_matrix}
\renewcommand{\arraystretch}{1.15}
\small
\begin{tabularx}{\textwidth}{|p{2.8cm}|p{2.5cm}|p{2.8cm}|X|}
\hline
\textbf{Threat} &
\textbf{STRIDE Category} &
\textbf{OWASP Mapping} &
\textbf{Researched Mitigations} \\
\hline
Direct Prompt Injection &
Tampering / Elevation of Privilege &
LLM01: Prompt Injection &
Dual-LLM/LMM verification architectures, Strict separation of data and instructions, Clear input boundaries \\
\hline
Indirect Prompt Injection &
Tampering &
LLM01: Prompt Injection &
Isolated sandbox parsing for untrusted data (web pages, PDFs), Semantic boundary sanitization before the LLMs receive an input \\
\hline
Cross-Modal Injection (Pixel/Audio Jailbreaks) &
Tampering / Spoofing &
LLM01: Prompt Injection (Multimodal Extension) &
OCR pre-scans on images, Adversarial perturbation filtering, Vision/audio encoder alignment checks \\
\hline
Jailbreaking / Guardrail Bypass &
Elevation of Privilege &
LLM01: Prompt Injection &
Layered system instructions, Independent guardrail models (e.g., Llama Guard), Defensive alignment tuning (RLHF/DPO) \\
\hline
System Prompt Leakage &
Information Disclosure &
LLM02: Sensitive Information Disclosure &
Strict prompt isolation, Programmatic output regex/semantic filtering, Structural design preventing system prompt disclosure \\
\hline
Multimodal Data Poisoning (Training/Fine-tuning) &
Tampering &
LLM03: Training Data Poisoning &
Cryptographic provenance tracking for training data and model weights, Robust data-cleaning pipelines, Gradient anomaly detection \\
\hline
\end{tabularx}
\end{table*}

\section{Score and Prioritize with DREAD}
Each threat is scored across five dimensions: Damage, Reproducibility, Exploitability, Affected Users, Discoverability.
\begin{table*}[ht]
\centering
\scriptsize
\renewcommand{\arraystretch}{1.15}
\caption{MITRE ATLAS Checklist - Threat Register}
\label{tab:mitre_atlas_threat_register}
\begin{tabularx}{\textwidth}{|X|X|X|X|X|X|}
\hline
\textbf{Threat} &
\textbf{STRIDE Component} &
\textbf{Asset Targeted} &
\textbf{Trust Boundary} &
\textbf{ATLAS Tactic} &
\textbf{ATLAS Technique} \\
\hline
Direct Prompt Injection &
Tampering &
System Prompt \& Intent Parser &
User $\to$ Orchestrator &
Initial Access / Execution &
AML.T0051: LLM Prompt Injection \\
\hline
Behavioral Jailbreaking \newline (e.g., Crescendo) &
Elevation of Privilege &
Safety Alignment Weights \& Policies &
User $\to$ LLM Core Engine &
Privilege Escalation / Defense Evasion &
AML.T0054: LLM Jailbreak \\
\hline
System Prompt Leakage &
Information Disclosure &
Hidden Architecture / Context Strings &
Core Engine $\to$ User &
Collection / Exfiltration &
AML.T0056: LLM Meta Prompt Extraction \\
\hline
Resource Exhaustion \newline (Token/Context Flooding) &
Denial of Service &
Compute Allocation \& Memory Buffers &
User $\to$ Gateway &
Impact &
AML.T0029: Denial of ML Service \\
\hline
Indirect RAG Injection \newline (Retrieval Poisoning) &
Tampering &
Vector Knowledge Base Embeddings &
External DB $\to$ Orchestrator &
Initial Access / Execution &
AML.T0051.001: Indirect Prompt Injection \\
\hline
Cross-Modal Exploitation \newline (Pixel-Hidden Injection) &
Spoofing / Tampering &
Vision/Audio Multimodal Encoders &
External Media Input $\to$ LMM &
ML Attack Staging / Initial Access &
AML.T0043: Craft Adversarial Data \\
\hline
\end{tabularx}
\end{table*}
\FloatBarrier

\section{Attack Taxonomy}
% TODO: taxonomy tree below is a working draft — cross-check final categories against the papers before finalizing. Every threat should map to exactly one category.

\begin{table*}[ht]
\centering
\caption{Attack Taxonomy for LLM/LMM Threats}
\label{tab:attack_taxonomy}
\renewcommand{\arraystretch}{1.15}
\small
\begin{tabularx}{\textwidth}{|p{2.6cm}|p{2.8cm}|p{3.2cm}|X|}
\hline
\textbf{Category} &
\textbf{STRIDE} &
\textbf{Example Threats} &
\textbf{Example Mitigations} \\
\hline
Prompt Injection &
Tampering, Elevation of Privilege &
Direct, indirect, cross-agent injection &
Prompt isolation, input validation, tool authorization \\
\hline
Adversarial Manipulation &
Tampering &
Data poisoning, retrieval poisoning, adversarial examples &
Data validation, provenance checks, retrieval filtering \\
\hline
Hallucination Exploitation &
Information Disclosure, integrity-related impacts &
Fabricated citations, false APIs, unsafe advice &
Grounding, confidence estimation, human review \\
\hline
Semantic Misalignment &
Elevation of Privilege, Tampering &
Goal manipulation, instruction conflicts &
Alignment testing, policy constraints, evaluation suites \\
\hline
Cross-modal Attacks &
Tampering &
OCR injection, hidden text, image-based prompts &
Content sanitization, modality-specific filtering \\
\hline
\end{tabularx}
\end{table*}
\FloatBarrier

\end{document}

